\chapter*{后记}
\phantomsection
\addcontentsline{toc}{chapter}{后记}
\markboth{后记}{}
\thispagestyle{plain}

\section{第一位读者来信}
\begin{myquote}
\textit{------------------第一章------------------}

% 我拜读了这一篇大作后，不禁大吃一惊、叹为观止。
我大吃一惊、瞠目结舌、叹为观止、拍案叫绝地拜读了这一篇空山啼竹素女愁、石破天惊惊秋雨的大作。

我觉得你写的内容就很有这类书该有的感觉，你真的很有写书的天赋啊（就是抛出新概念，然后列举生动的事例进行解释的写作手法）！

而且你的阅读量和词汇量远比我想象中要丰富啊——包括月光宝盒、吸星大法、一剪梅、基督山伯爵……

你确实做到了把专业话题讲得通俗易懂。其实我之前一直不是很了解熵（虽然每一本写信息的书基本上都要提一嘴，比如《数字化生存》，但它就没有怎么解释，对我们文科生很不友好），但现在我确实看懂了。它是混乱的，但信息世界会让它变得更简洁。

而且我觉得那个“先有果后有因”的事例讲得好有趣，就是“会先被打死再看到要被打死”。

还有就是你的思考真的挺有深度的，新鲜概念很多，其实我不喜欢看这种书的一个原因就是每本书的概念都大差不差，像老派爱情电影一样不断重复，但是你就能提出了很多新概念（或许是因为时代在进步）。

你写得很有创新性的一个要点就是把这本书写得很像教科书，有很有意思的思考题，我觉得很新颖、很有趣。

关于共识的那一处内容，我看到你用了meme（梗），我之前看自私的基因的时候，它一直被翻译为“模因”。这方面确实是你讲得更明白些，不是那么老学究……


\textit{------------------第二章------------------}

你肯定有大文学家的天赋！

第二章里面我最喜欢的是双重视角的内容（炼油厂和油田）——不同视角的转换就很清晰且全面地写出了过多数据的危害（而且其中举例Apple Watch的那一段我很喜欢，我觉得排版很好看，长短句的结合很有节奏感，还有一种从之前的原理到这里的事例的“抽离感”）。

之前是谁说自己语文不好的？我觉得你的语文水平蛮好的（甚至营造出了戏剧化小说的感觉）！

而且你的历史知识挺丰富的。我看到有一段的比喻是圈地运动，我觉得你历史学得挺好的，我读了三四遍才想通这个比喻里具体的对照点（也可能是我英国史学得不好的原因）。这里本身就有一种对资本家的批判吧？或者可以理解为本质一样。

我发现你们这种AI行业的人真的很喜欢典故之类的内容（都在说巴别塔、香农等等），这是因为你们开了什么信息史研究课，要完成才给分配工作吗？

但作为一个讨厌概率的高中生（期望算不清楚、马上考T8版），数学题我是全跳过的（私密马赛bushi），就是说数学题是“何意为”啊……

好的，我上面的玩梗“停之停之”，总之我完全能看懂你想表达什么。作为读者（反正姐也不想进业内），我觉得完全足够了，这本书非常优秀！

\textit{------------------第七章------------------}

这种压缩产生了副产品：当模型被问到“中国的首都是哪里”时，它并不须要真的“知道”这个事实，它只须要根据训练数据中的这个问题与“北京”共现的统计规律，给出高概率的回答——我想我已经理解LLM回答问题的原理了。

“思想实验”这个词好有美式风格啊！你现在真的变成一个美式男孩了啊！

你是希望未来的LLM能够将记忆引入它的思维过程中吗——是希望它能思考吗（但是记忆是人脑的功能，我们现在哲学的观点是只有人脑才能思考。我们认为猪脑、人工智能等都无法思考）？

我感觉你确实是希望它能变成人类呢。

\hfill ——姜彦孜
\end{myquote}




\newpage
\section*{高一片叶子}

\begin{myquote}
我不要求你们比其他人领先很多，只要求你们比其他人高出“一片叶子”。只要做到这一点，天空中所有的阳光和雨露都是你的。

\hfill ——曹家树
\end{myquote}

\vspace{1em}

彦孜：

你说“信息世界会让熵变得更简洁”。这句话写得比我整本书都好。

数学全跳过了，没事。序章有个脚注你大概也跳过了。AI能让能源变成算力，算力变成劳动力。以后能源和算力就是最底层的生产力，算力的单位叫Token。我们会需要更多的能源，算力会和电一样普及，变成比黄金还要重要、还要有用的东西。

\vspace{1em}

大家焦虑AI，很多时候是把“能干活”和“有用”画了等号。
你知道我最近老在想什么吗？雅典。你历史学的比我好，肯定比我熟。奴隶干活，自由民搞哲学、搞戏剧、搞政治。

\vspace{1em}

AI要是把活都干了，我们是不是全变自由民了？
但雅典的自由民之间，照样分三六九等，照样争夺执政官的位置，照样打官司，照样在奥运会上拼到吐血。体力劳动交出去了，人跟人之间的事一件没少。

\vspace{1em}

这和有没有AI没有关系。

退一万步讲，就算某一天AI取代了我们所有人的工作，只要社会还存在，只要政府还能对社会财富进行二次分配，我们的生活和过去可能也不会有太大的不同。

我们不需要和AI竞争，因为它比我们所有人都更能干活。生活在我们周围的，仍然会是过去那些亲人、朋友、同学、同事和陌生人。

\vspace{1em}

你问我是不是希望AI变成人类——不是。
我不希望它变成人类。
我希望它始终是工具，是足够聪明、足够强大的工具。

\newpage

但就算有一天，它真的学会了和人类一样“思考”，它也仍然替代不了人与人之间那种真实的关系感。

那个情书的故事，重点不是AI写得好不好，重点是你不知道对方爱不爱你。得自己去确认，去争取，冒着被拒绝的风险写那些笨句子。当前，LLM被设计为天生谄媚的工具，是永远说“好的，主人”的魔镜，它不会拒绝你。让关系有重量的，恰恰是不确定，是等待，是误解，是试探，是你坚定不退缩的选择，是明知道可能没有回音，还是把心意递出去。

\vspace{1em}

高一分文理科的时候，你说你想选文科，“不想进业内”，老实说，我有点意外，也有点失望。
还有三个月你就要高考了，选什么学校，读什么专业，坦率地说，在你跟我说想选文科的时候，我就知道在这件事上，我给不了你什么指导。

但我觉得这样也很好，就按照你自己的想法来。尺子不是丈量龙的，是召唤龙的。你拿什么标准衡量自己，就会长成什么样子。


未来的社会需要与众不同的人。你和我玩的那些梗、那些跳脱的诗句、那些“私密马赛bushi”、那些让你成为你而不是别人的东西，AI学不走，也复制不了。找一个你愿意较真的地方就行。你有你自己的Token，注意力就是你的Token。花在哪里，比有多少更重要。不用赢，别投降。

\vspace{1em}

四部篇章，六道魔法，几十个概念。写的时候，我想把所有我知道的都写给你，每一个我都觉得很重要，每一个我都不舍得删。但如果这本书真的能对你有用，你仍然要从书里找出你最在意的那一两句话，找到属于你的“稍纵即逝的八角分”。
如果只能留一个信号，我想留这句——AI能干所有分布之内的事，它最擅长的就是平均。

\vspace{1em}

\textbf{而你能给这个世界的，在分布之外。}

高出一片叶子就够了，你自己的那一片。

\vspace{2em}
\begin{flushright}
爸爸
\end{flushright}
